\section{模型评价、局限与改进方向}\label{sec:conclusion}

本节不再复述各问结论（逐条答案与出处见第 1.5 节表~\ref{tab:v42-answer-index}，各问的数值结果与适用范围见相应小节），只评价模型本身的结构优缺点、明确解释边界，并给出后续改进的优先级. 

\subsection{模型的主要优势}
本文的主要特点在于解析结构、决策目标与数据证据之间保持一致，并通过可证明的数学性质和可复核的数值证据形成完整的方法链条.
\setcounter{obsitem}{0}
\obsitem{明确了参数的可识别范围与分析边界：}
在固定域质量下，配比系数与质量交互系数存在严格的线性相关关系（命题~\ref{prop:q1-quality-collinearity-v15}）；锚点数据无法单独识别跨来源配比迁移系数（命题~\ref{prop:q2-lambda-unidentified-v15}）；对于未映射域，完整质量由命题（命题~\ref{prop:background-quality-bound}） 给出可识别区间. 这些结果将参数识别条件直接写入模型结构，使后续各问的推断范围具有明确的数学依据.
\obsitem{能够定量刻画预测误差对决策结果的影响：}
候选集合上的一致误差界通过三次比较传递为选择后悔界（命题~\ref{prop:q1-decision-regret-v37}）；有序分位规划的全局目标间隙则通过强凸性进一步转化为共享斜率误差界（推论~\ref{cor:q4-slope-gap-v37}）. 因此，模型不仅给出预测值和最优选择，还能够对预测精度、优化精度与最终决策质量之间的关系进行定量刻画，使“预测是否准确”和“决策是否可靠”处于统一的误差分析框架中.
\obsitem{预算变化下的最优配置具有严谨的数学证明：}
仿射下包络的凹性与严格单调性保证全局最优参数量随预算严格增加，同时最优质量不下降（定理~\ref{thm:q3-global-monotonicity-v37}）. 进一步地，有限预算下的质量饱和与质量成本割线比有界之间具有充要关系（定理~\ref{thm:q3-finite-saturation-iff-v37}）. 由此，参数规模、数据规模与质量投入随预算变化的规律不仅由数值结果展示，还具有相应的结构性判据.
\obsitem{主要数值结果均具有可复核的求解精度保证：}
45 个名义配置的最大目标间隙、预算相对残差、稳健配置的基准误差上界以及 279 份有序分位拟合的目标间隙均被统一记录并归档. 这些数值证书与前述解析结果相互配合，使模型从理论推导、数值求解到最终决策均具有明确的精度依据，并提高了整个建模流程的可复核性与一致性.


\subsection{模型的主要局限}
\setcounter{obsitem}{0}
\obsitem{质量衡量仍是间接的:}
本文的质量变量 $Q$ 由文本信号构造, 并非直接的人类评价或真实下游训练效果. 六个可映射域覆盖锚点配比的 $31.94\%$, 因此质量分析主要反映当前指标体系下的相对差异.
\obsitem{部分结果依赖模型设定:}
第二、三问中的质量响应主要来自经真实数据校准的半合成数据, 成本函数也按照题目给定的训练、质量处理和注意力开销构造. 因此, 得到的最优配置与资源路径对应本文所采用的质量响应和成本设定.
\obsitem{未来能力验证数据还比较少:}
第四问的严格前沿样本为 33 条, 开发期包含 16 个预测标签. 现有数据用于构建条件能力前沿并开展滚动检验和整组织留出分析, 后续随着更多模型和长期数据积累, 可进一步验证 12/24 月能力情景.

\subsection{模型推广}
本文的三层结构可迁移到其他"有限预算下的多要素配置"问题: 先用部分识别界定数据能支持的质量与配比信息, 再把不可直接观测的要素折算为有效资源量, 最后用成本下包络把多维配置问题约化为可全局比较的一维尺度问题. 

其中第二、三问依赖"质量响应单调且存在有效数据折算"这一结构, 因此在语料质量存在非单调效应或多模态训练等场景下, 需要先重新检验该前提.
\subsection{AI 使用声明}
本文章使用了chatGPT工具，版本为GPT-6 Sol / GPT-6 Luna开发公司为 OpenAI版本颁布日期2026年9月22日. \\
主要用于辅助概念理解（\url{https://chatgpt.com/share/6ab68617-b2b0-83e8-b719-ff1822b2f894}）、\\
资料搜索与整理（\url{https://chatgpt.com/share/6ab6857e-54fc-83ee-ac02-d0bb01562831}）、程序实现与调试、数值结果核对、文字整理及版面核查（\url{https://chatgpt.com/share/6ab68617-b2b0-83e8-b719-ff1822b2f894}）\\（\url{https://chatgpt.com/s/cx_6ab68696564c8191a4953b0863f7668e}）\\
全部模型设定、数学推导、结果解释与结论由参赛作者完成并逐项复核.\\
AI 未用于替代核心建模、推导与论证. 参赛作者对提交稿的模型设定、数据来源、计算结果和结论承担最终复核责任. 本文已列明的程序输出与证据边界供逐项复核.